%% ========================================================================
%% Chapter 7: Bash Shell
%% ========================================================================

\chapter{Bash Shell}
\label{ch:bash-shell-basic}

\chapterhero{cyan}{Bangun kebiasaan kerja terminal yang efisien lewat konfigurasi shell, PATH, alias, fungsi, dan teknik interaksi Bash yang lebih produktif.}{shell}{PATH}{alias}

Setiap kali kamu membuka terminal di Linux, ada program yang langsung
aktif menunggu perintahmu. Program itu bukan aplikasi biasa, melainkan
\textit{shell}, yaitu penerjemah antara perintah yang kamu ketik dan
sistem operasi yang menjalankannya. Shell membaca baris teks yang kamu
tulis, memprosesnya, lalu menyuruh kernel Linux melakukan sesuatu.

Bash, singkatan dari \textit{Bourne Again Shell}, adalah shell paling
umum di Linux. Hampir semua distribusi Linux menjadikan Bash sebagai
shell default. Nama ``Bourne Again'' adalah permainan kata: Bash
merupakan penerus \cmd{sh} (Bourne Shell) yang lahir di era Unix tahun
1979. Selain Bash, ada shell lain seperti \cmd{zsh} dan \cmd{fish},
tetapi Bash tetap menjadi pilihan utama untuk administrasi sistem dan
penulisan skrip karena tersedia hampir di semua sistem Linux.

Bayangkan Bash seperti dapur di sebuah restoran. Kamu sebagai
administrator adalah kepala koki. Sistem operasi adalah dapur beserta
semua peralatannya. Bash adalah commis chef yang menerima instruksimu,
mengambil bahan dari lemari yang tepat (direktori), menggunakan alat
yang sesuai (program), dan menyajikan hasilnya. Tanpa shell, kamu harus
berbicara langsung dengan kernel dalam bahasa yang jauh lebih rendah.
Dengan Bash, kamu cukup mengetik perintah dalam bahasa manusia yang
terstruktur.

\textbf{Yang Akan Kamu Pelajari}

Setelah menyelesaikan bab ini, kamu akan mampu:

\begin{enumerate}
    \item menjelaskan perbedaan antara shell login dan shell interaktif serta
          mengonfigurasi \file{~/.bashrc} dan \file{~/.bash\_profile};
    \item mengelola variabel lingkungan dan memanipulasi \var{PATH} untuk
          menjalankan perintah dari direktori pribadi;
    \item membuat alias dan fungsi shell untuk mempercepat pekerjaan
          administrasi sehari-hari;
    \item menggunakan wildcard, brace expansion, dan fitur history Bash
          secara efektif dan aman;
    \item menerapkan quoting dan escaping yang benar untuk menangani
          nama file dengan spasi dan karakter khusus.
\end{enumerate}

\begin{notebox}
Seluruh praktek pada bab ini dijalankan pada direktori
\dir{~/lab-os/chapter06-bash/}.
Buat direktori ini sebelum memulai:
\begin{lstlisting}[language=bash]
mkdir -p ~/lab-os/chapter06-bash
cd ~/lab-os/chapter06-bash
\end{lstlisting}
\end{notebox}

%% ------------------------------------------------------------------------
\section{Pengenalan Bash dan Konfigurasi}
\label{sec:pengenalan-bash-konfigurasi}

Bash bukan sekadar tempat mengetik perintah. Ia memiliki sistem
konfigurasi yang memungkinkan kamu mengatur tampilan prompt, mendefinisikan
variabel, dan menjalankan perintah tertentu setiap kali shell dibuka.
Konfigurasi ini disimpan dalam file teks biasa di direktori home.

Ada dua konteks berbeda di mana Bash bisa berjalan. \textit{Shell login}
adalah shell pertama yang muncul saat kamu masuk ke sistem, misalnya
lewat SSH atau login teks. Shell ini membaca \file{~/.bash\_profile} (atau
\file{~/.profile} jika \file{~/.bash\_profile} tidak ada). \textit{Shell
interaktif non-login} adalah shell yang muncul saat kamu membuka jendela
terminal baru di dalam sesi grafis yang sudah aktif. Shell ini membaca
\file{~/.bashrc}.

Praktik standar yang umum dilakukan adalah membuat \file{~/.bash\_profile}
memanggil \file{~/.bashrc}, sehingga semua konfigurasi hanya perlu ditulis
di satu tempat. Dengan cara ini, alias, fungsi, dan variabel yang kamu
definisikan di \file{~/.bashrc} akan aktif baik di shell login maupun
shell interaktif.

Prompt Bash dikontrol oleh variabel \var{PS1}. Nilai default-nya biasanya
menampilkan nama pengguna, hostname, dan direktori aktif dalam format
\texttt{user@host:dir\$}. Kamu bisa mengubahnya untuk menampilkan informasi
tambahan. Beberapa kode khusus yang berguna di dalam \var{PS1}:
\texttt{\textbackslash u} untuk nama pengguna, \texttt{\textbackslash h}
untuk hostname pendek, \texttt{\textbackslash w} untuk path direktori aktif,
dan \texttt{\textbackslash t} untuk waktu dalam format HH:MM:SS.

Perintah \cmd{source} atau ekuivalennya \cmd{.} (titik) digunakan untuk
memuat ulang file konfigurasi tanpa perlu menutup dan membuka terminal.
Setelah mengedit \file{~/.bashrc}, jalankan \cmd{source ~/.bashrc} agar
perubahan langsung aktif di sesi yang sedang berjalan.

\subsection*{Praktek 6.1: Modifikasi Prompt dan Konfigurasi Shell}

Tujuan praktek ini adalah memahami peran file konfigurasi Bash dan
mengubah tampilan prompt secara aman.

\begin{enumerate}
    \item Periksa shell yang sedang aktif dan versi Bash yang terpasang:
\begin{lstlisting}[language=bash, caption={Memeriksa shell aktif}]
echo "Shell login : $SHELL"
echo "Shell aktif : $0"
bash --version | head -n 1
\end{lstlisting}
    \textit{Amati:} \var{SHELL} menunjukkan shell default pengguna yang
    terdaftar di sistem, sementara \cmd{\$0} menunjukkan nama proses shell
    yang sedang berjalan saat ini. Keduanya bisa berbeda.

    \item Lihat isi \file{~/.bashrc} yang sudah ada:
\begin{lstlisting}[language=bash, caption={Melihat konfigurasi .bashrc saat ini}]
ls -la ~ | grep -E 'bashrc|bash_profile|profile'
wc -l ~/.bashrc
\end{lstlisting}

    \item Buat backup \file{~/.bashrc} sebelum mengubah apapun:
\begin{lstlisting}[language=bash, caption={Membuat backup .bashrc}]
cp ~/.bashrc ~/.bashrc.bak-chapter06
echo "Backup selesai: ~/.bashrc.bak-chapter06"
\end{lstlisting}

    \item Tambahkan konfigurasi prompt kustom ke \file{~/.bashrc}. Blok ini
          ditandai komentar agar mudah diidentifikasi dan dihapus jika perlu:
\begin{lstlisting}[language=bash, caption={Menambahkan konfigurasi prompt ke .bashrc}]
nano ~/.bashrc
# Tambahkan blok berikut di akhir file

# === Konfigurasi Bab 6 Bash Shell ===
# Prompt: user@host:dir [waktu]$
PS1='\u@\h:\w [\t]\$ '
export EDITOR=nano
export VISUAL=nano
# === Akhir Konfigurasi Bab 6 ===

\end{lstlisting}

    \item Terapkan konfigurasi tanpa menutup terminal:
\begin{lstlisting}[language=bash, caption={Memuat ulang konfigurasi Bash}]
source ~/.bashrc
echo "Prompt aktif: $PS1"
echo "Editor: $EDITOR"
\end{lstlisting}
    \textit{Amati:} prompt terminal berubah menampilkan waktu di dalam
    tanda kurung siku. Setiap kali kamu menekan Enter, waktu di prompt
    diperbarui.

    \item Periksa apakah \file{~/.bash\_profile} sudah memanggil
          \file{~/.bashrc}. Jika belum, buat konfigurasikan:
\begin{lstlisting}[language=bash, caption={Memastikan .bash\_profile memanggil .bashrc}]
if ! grep -q '\.bashrc' ~/.bash_profile 2>/dev/null; then
    nano ~/.bash_profile
# Tambahkan blok berikut bila belum ada

# Panggil .bashrc dari shell login
if [ -f ~/.bashrc ]; then
    . ~/.bashrc
fi
    echo ".bash_profile diperbarui"
else
    echo ".bash_profile sudah memanggil .bashrc"
fi
\end{lstlisting}
\end{enumerate}

\begin{challengebox}
Tambahkan informasi direktori aktif dan waktu sekarang ke prompt secara
lebih informatif. Ubah \var{PS1} sehingga direktori aktif ditampilkan
dengan warna biru dan waktu dengan warna hijau menggunakan kode warna ANSI:
\texttt{\textbackslash[\textbackslash 033[34m\textbackslash]} untuk biru
dan \texttt{\textbackslash[\textbackslash 033[32m\textbackslash]} untuk
hijau. Gunakan perintah \cmd{echo} dari Bab~2 untuk menguji apakah
terminal mendukung warna dengan \cmd{echo -e \textbackslash
033[32mHijau\textbackslash033[0m}. Pastikan prompt kembali ke warna normal
dengan kode reset \texttt{\textbackslash[\textbackslash033[0m\textbackslash]}
di akhir.
\end{challengebox}

%% ------------------------------------------------------------------------
\section{Variabel Lingkungan dan PATH}
\label{sec:variabel-lingkungan-path}

Variabel di Bash bekerja seperti label pada kotak penyimpanan. Kamu
memberi nama pada sebuah kotak, mengisi nilai ke dalamnya, dan
mengambilnya kembali saat dibutuhkan. Ada dua jenis variabel di Bash.
\textit{Variabel shell lokal} hanya hidup di dalam shell saat ini dan
tidak diwariskan ke proses anak. \textit{Variabel lingkungan} (environment
variable) ditandai dengan perintah \cmd{export} dan diwariskan ke setiap
program yang dijalankan dari shell tersebut.

Sintaks dasar: \texttt{NAMA=nilai} membuat variabel lokal. Tidak boleh ada
spasi di sekitar tanda sama dengan. Untuk membacanya, gunakan
\texttt{\$NAMA} atau \texttt{\$\{NAMA\}}. Kurung kurawal berguna saat
variabel disambung langsung dengan teks, misalnya
\texttt{\$\{NAMA\}\_backup}.

Variabel \var{PATH} adalah daftar direktori yang dipisahkan tanda titik
dua. Saat kamu mengetik sebuah perintah tanpa path lengkap, Bash mencari
program dengan nama itu di setiap direktori dalam \var{PATH} dari kiri ke
kanan. Misalnya, nilai \var{PATH} seperti
\texttt{/usr/local/bin:/usr/bin:/bin} berarti Bash mencari di
\dir{/usr/local/bin} dulu, kemudian \dir{/usr/bin}, lalu \dir{/bin}.

Menambahkan direktori pribadi ke \var{PATH} sangat berguna. Kamu bisa
menyimpan skrip-skrip yang sering dipakai di \dir{~/bin} dan memanggilnya
dari mana saja tanpa perlu mengetik path lengkap. Caranya dengan menambahkan
\texttt{export PATH="\$HOME/bin:\$PATH"} ke \file{~/.bashrc}. Perhatikan
\texttt{\$PATH} di akhir: ini penting agar direktori sistem yang sudah ada
tidak hilang.

Beberapa variabel lingkungan bawaan yang penting: \var{HOME} menyimpan path
home directory, \var{USER} menyimpan nama pengguna yang sedang login,
\var{SHELL} menyimpan path shell default, \var{PWD} menyimpan direktori
aktif saat ini, dan \var{OLDPWD} menyimpan direktori sebelumnya sehingga
\cmd{cd -} bisa kembali ke sana.

\subsection*{Praktek 6.2: Kelola Variabel Lingkungan}

Tujuan praktek ini adalah memahami perbedaan variabel lokal dan
variabel lingkungan, serta menambahkan direktori pribadi ke \var{PATH}.

\begin{enumerate}
    \item Buat variabel lokal dan cek apakah ia diwariskan ke subshell:
\begin{lstlisting}[language=bash, caption={Variabel lokal tidak diwariskan ke subshell}]
PROYEK="os-book"
echo "Di shell ini: $PROYEK"

bash -c 'echo "Di subshell: $PROYEK"'
\end{lstlisting}
    \textit{Amati:} baris kedua menampilkan string kosong karena variabel
    lokal tidak diwariskan ke proses anak.

    \item Ubah menjadi variabel lingkungan dengan \cmd{export}:
\begin{lstlisting}[language=bash, caption={Variabel lingkungan diwariskan ke subshell}]
export PROYEK="os-book"
bash -c 'echo "Di subshell: $PROYEK"'
\end{lstlisting}
    \textit{Amati:} kali ini subshell berhasil membaca nilai variabel.

    \item Lihat semua variabel lingkungan yang aktif saat ini:
\begin{lstlisting}[language=bash, caption={Melihat variabel lingkungan}]
printenv | grep -E '^(HOME|USER|SHELL|PATH|PROYEK)='
\end{lstlisting}

    \item Lihat isi \var{PATH} saat ini dan cari lokasi beberapa perintah:
\begin{lstlisting}[language=bash, caption={Melihat PATH dan lokasi perintah}]
echo "$PATH" | tr ':' '\n'
which bash
which ls
type ls
\end{lstlisting}
    \textit{Amati:} perintah \cmd{which} menampilkan path file eksekutabel,
    sedangkan \cmd{type} bisa menampilkan bahwa \cmd{ls} adalah alias
    jika alias sudah didefinisikan di sistem.

    \item Buat direktori \dir{~/bin} dan tambahkan ke \var{PATH} secara
          permanen di \file{~/.bashrc}:
\begin{lstlisting}[language=bash, caption={Menambahkan ~/bin ke PATH}]
mkdir -p ~/bin

nano ~/.bashrc
# Tambahkan blok berikut di akhir file

# Tambahkan ~/bin ke PATH jika belum ada
if [[ ":$PATH:" != *":$HOME/bin:"* ]]; then
    export PATH="$HOME/bin:$PATH"
fi

source ~/.bashrc
echo "$PATH" | tr ':' '\n' | head -5
\end{lstlisting}

    \item Buat skrip sederhana di \dir{~/bin} dan jalankan dari direktori
          yang berbeda:
\begin{lstlisting}[language=bash, caption={Membuat dan menguji skrip di ~/bin}]
nano ~/bin/sysinfo
# Tulis isi berkas berikut
#!/bin/bash
echo "User   : $(whoami)"
echo "Host   : $(hostname)"
echo "Uptime : $(uptime -p)"
echo "Disk / : $(df -h / | awk 'NR==2 {print $5}') terpakai"

chmod +x ~/bin/sysinfo

cd /tmp
sysinfo
cd ~/lab-os/chapter06-bash
\end{lstlisting}
    \textit{Amati:} skrip \cmd{sysinfo} bisa dipanggil dari \dir{/tmp}
    tanpa menuliskan path lengkap, karena \dir{~/bin} sudah masuk ke
    \var{PATH}.
\end{enumerate}

\begin{challengebox}
Tambahkan direktori \dir{~/lab-os/chapter06-bash/bin} ke \var{PATH} secara
permanen di \file{~/.bashrc} menggunakan pola pengecekan yang sama pada
langkah 5. Kemudian buat skrip \cmd{cuaca-disk} di direktori tersebut yang
menampilkan status setiap filesystem yang terpasang: jika pemakaian di atas
80\% tampilkan peringatan ``PENUH'', jika di bawah 50\% tampilkan ``AMAN'',
dan di antaranya tampilkan ``NORMAL''. Gunakan \cmd{df} dan \cmd{awk}
untuk mengekstrak persentase. Verifikasi skrip bisa dipanggil dari
direktori \dir{/tmp} tanpa path lengkap.
\end{challengebox}

%% ------------------------------------------------------------------------
\section{Alias dan Fungsi Shell}
\label{sec:alias-fungsi-shell}

Alias adalah nama pendek untuk perintah panjang yang sering kamu gunakan.
Bayangkan alias seperti kontak di ponsel: daripada menghafal dan mengetik
nomor panjang setiap kali, kamu cukup menekan satu nama. Alias cocok untuk
perintah yang sederhana dan tidak memerlukan logika.

Sintaks alias: \texttt{alias nama='perintah panjang'}. Untuk melihat semua
alias yang aktif, ketik \cmd{alias} tanpa argumen. Untuk menghapus alias,
gunakan \cmd{unalias nama}. Alias yang didefinisikan di terminal hanya
bertahan selama sesi itu. Agar permanen, tambahkan ke \file{~/.bashrc}.

Fungsi shell adalah langkah berikutnya. Fungsi bisa menerima argumen,
menjalankan beberapa perintah, menggunakan variabel lokal, dan mengembalikan
nilai. Gunakan fungsi saat logika yang kamu butuhkan lebih dari sekadar
singkatan perintah. Dua cara mendefinisikan fungsi yang keduanya valid:

\begin{lstlisting}[language=bash, caption={Dua cara mendefinisikan fungsi shell}]
# Cara 1: sintaks tradisional
nama_fungsi() {
    perintah
}

# Cara 2: dengan kata kunci function
function nama_fungsi {
    perintah
}
\end{lstlisting}

Di dalam fungsi, \texttt{\$1}, \texttt{\$2}, dan seterusnya adalah
argumen yang diberikan saat fungsi dipanggil. Gunakan \cmd{local} untuk
mendeklarasikan variabel yang hanya hidup di dalam fungsi. Tanpa \cmd{local},
variabel di dalam fungsi akan mengubah variabel global dengan nama yang sama.

Perintah \cmd{type} sangat berguna untuk memverifikasi apakah sebuah nama
adalah alias, fungsi, perintah builtin Bash, atau program eksternal. Cukup
ketik \texttt{type nama} untuk melihat jenisnya.

\subsection*{Praktek 6.3: Buat Alias dan Fungsi Produktivitas}

Tujuan praktek ini adalah membuat alias dan fungsi yang berguna untuk
pekerjaan sehari-hari di terminal.

\begin{enumerate}
    \item Tambahkan beberapa alias produktivitas ke \file{~/.bashrc}:
\begin{lstlisting}[language=bash, caption={Menambahkan alias ke .bashrc}]
nano ~/.bashrc
# Tambahkan blok berikut di akhir file

# === Alias Bab 6 ===
alias ll='ls -lah --color=auto'
alias la='ls -A --color=auto'
alias hist10='history | tail -n 10'
alias grep='grep --color=auto'
alias ..='cd ..'
alias ...='cd ../..'
alias chapter06='cd ~/lab-os/chapter06-bash'
# === Akhir Alias Bab 6 ===


source ~/.bashrc
\end{lstlisting}

    \item Uji alias yang baru dibuat:
\begin{lstlisting}[language=bash, caption={Menguji alias}]
ll ~/lab-os/chapter06-bash
type ll
type hist10
type ..
\end{lstlisting}
    \textit{Amati:} perintah \cmd{type} menampilkan bahwa \cmd{ll} adalah
    alias beserta perintah asli di baliknya.

    \item Tambahkan fungsi \cmd{backup\_conf} yang membuat backup file
          konfigurasi dengan timestamp:
\begin{lstlisting}[language=bash, caption={Menambahkan fungsi backup\_conf ke .bashrc}]
nano ~/.bashrc
# Tambahkan blok berikut di akhir file

# === Fungsi Bab 6 ===
backup_conf() {
    if [ $# -ne 1 ]; then
        echo "Penggunaan: backup_conf <file>"
        return 1
    fi
    local src="$1"
    local dst="$HOME/lab-os/chapter06-bash/backup"
    if [ ! -f "$src" ]; then
        echo "File tidak ditemukan: $src"
        return 2
    fi
    mkdir -p "$dst"
    local nama_backup
    nama_backup="$(basename "$src").$(date +%F-%H%M%S).bak"
    cp -- "$src" "$dst/$nama_backup"
    echo "Backup disimpan: $dst/$nama_backup"
}
# === Akhir Fungsi Bab 6 ===


source ~/.bashrc
\end{lstlisting}

    \item Uji fungsi dengan membuat file konfigurasi contoh lalu
          mem-backup-nya:
\begin{lstlisting}[language=bash, caption={Menguji fungsi backup\_conf}]
mkdir -p ~/lab-os/chapter06-bash/backup
echo "PORT=8080\nHOST=localhost" > ~/lab-os/chapter06-bash/app.conf

backup_conf ~/lab-os/chapter06-bash/app.conf
backup_conf ~/lab-os/chapter06-bash/app.conf

ls -lah ~/lab-os/chapter06-bash/backup/
type backup_conf
\end{lstlisting}
    \textit{Amati:} setiap panggilan \cmd{backup\_conf} menghasilkan file
    backup dengan nama yang berbeda karena timestamp berbeda, sehingga tidak
    ada backup yang tertimpa.
\end{enumerate}

\begin{challengebox}
Buat fungsi \cmd{mkcd()} yang membuat direktori sekaligus berpindah ke
dalamnya dalam satu langkah. Fungsi ini harus menangani kasus berikut:
jika tidak ada argumen, tampilkan pesan penggunaan dan kembalikan kode
error 1; jika direktori sudah ada, pindah ke dalamnya saja tanpa error;
jika direktori berhasil dibuat, tampilkan pesan konfirmasi. Tambahkan
fungsi ini ke \file{~/.bashrc}, muat ulang, lalu uji dengan membuat
beberapa direktori bersarang sekaligus menggunakan opsi \texttt{-p} di
dalam fungsi.
\end{challengebox}

%% ------------------------------------------------------------------------
\section{Wildcards, Ekspansi, dan History}
\label{sec:wildcards-ekspansi-history}

Bash memiliki kemampuan untuk mengembangkan pola sebelum menjalankan
perintah. Proses ini disebut \textit{ekspansi}. Ada beberapa jenis
ekspansi yang penting diketahui.

\textit{Wildcard} atau globbing adalah pola untuk mencocokkan nama file.
Karakter \texttt{*} cocok dengan nol atau lebih karakter apa pun kecuali
garis miring. Misalnya, \texttt{*.log} cocok dengan semua file yang
berakhiran \texttt{.log}. Karakter \texttt{?} cocok dengan tepat satu
karakter. Misalnya, \texttt{app-0?.log} cocok dengan \texttt{app-01.log}
dan \texttt{app-02.log} tetapi tidak dengan \texttt{app-001.log}. Karakter
\texttt{[...]} mendefinisikan kelas karakter: \texttt{[abc]} cocok dengan
\texttt{a}, \texttt{b}, atau \texttt{c}, sedangkan \texttt{[0-9]} cocok
dengan satu digit.

\textit{Brace expansion} adalah ekspansi yang tidak bergantung pada nama
file yang ada. Pola \texttt{\{a,b,c\}} menghasilkan tiga kata terpisah.
Pola \texttt{\{01..05\}} menghasilkan urutan \texttt{01}, \texttt{02},
\texttt{03}, \texttt{04}, \texttt{05}. Brace expansion sangat efisien
untuk membuat banyak file atau direktori sekaligus.

\textit{Command substitution} mengganti ekspresi \texttt{\$(perintah)}
dengan output dari perintah tersebut. Ini berguna untuk menyimpan output
perintah ke variabel atau menggunakannya langsung dalam string.

History Bash menyimpan setiap perintah yang pernah kamu jalankan.
Tekan tombol panah atas untuk menelusuri perintah sebelumnya. Ketik
\cmd{history} untuk melihat daftar lengkap dengan nomor urut. Gunakan
\texttt{!n} untuk menjalankan ulang perintah nomor n, dan \texttt{!!}
untuk mengulang perintah terakhir. Pencarian interaktif dengan
\keystroke{Ctrl+R} adalah cara tercepat: tekan \keystroke{Ctrl+R} lalu
ketik sebagian perintah yang kamu cari.

\subsection*{Praktek 6.4: Gunakan Ekspansi dan History}

Tujuan praktek ini adalah menguasai wildcard, brace expansion, dan
fitur history untuk bekerja lebih cepat di terminal.

\begin{enumerate}
    \item Buat struktur file sampel menggunakan brace expansion:
\begin{lstlisting}[language=bash, caption={Membuat file sampel dengan brace expansion}]
cd ~/lab-os/chapter06-bash
mkdir -p sampel logs arsip

touch sampel/catatan-{a,b,c,d}.txt
touch logs/app-{01..05}.log
touch logs/error-{01..03}.log

ls sampel/
ls logs/
\end{lstlisting}
    \textit{Amati:} dengan satu perintah \cmd{touch}, kamu membuat banyak
    file sekaligus. Brace expansion berjalan sebelum perintah dieksekusi.

    \item Coba berbagai pola wildcard:
\begin{lstlisting}[language=bash, caption={Menggunakan wildcard untuk memfilter file}]
ls logs/*.log
ls logs/app-0?.log
ls logs/app-0[123].log
ls sampel/catatan-[a-c].txt
\end{lstlisting}

    \item Sebelum menjalankan perintah yang berdampak, preview dulu hasilnya
          dengan \cmd{echo}:
\begin{lstlisting}[language=bash, caption={Preview wildcard sebelum operasi nyata}]
echo logs/*.log
echo logs/error-*.log
\end{lstlisting}

    \item Pindahkan file log ke direktori arsip menggunakan wildcard yang
          sudah diverifikasi:
\begin{lstlisting}[language=bash, caption={Memindahkan file log ke arsip}]
mv logs/error-*.log arsip/
ls logs/
ls arsip/
\end{lstlisting}

    \item Jalankan beberapa perintah lalu gunakan fitur history:
\begin{lstlisting}[language=bash, caption={Menggunakan history Bash}]
pwd
ls -lah
date
whoami
history | tail -n 10
\end{lstlisting}

    \item Cari perintah tertentu dari history menggunakan \cmd{grep}:
\begin{lstlisting}[language=bash, caption={Mencari perintah di history}]
history | grep "ls logs"
\end{lstlisting}
    \textit{Amati:} kamu bisa menjalankan ulang perintah yang ditemukan
    dengan mengetik \texttt{!} diikuti nomor historynya. Coba \texttt{!}
    diikuti nomor salah satu perintah di atas.
\end{enumerate}

\begin{challengebox}
Gunakan brace expansion untuk membuat 5 direktori sekaligus dengan nama
\texttt{bab-\{01..05\}} di dalam \dir{~/lab-os/chapter06-bash/proyek}.
Di dalam setiap direktori, buat tiga file: \texttt{rencana.txt},
\texttt{progres.txt}, dan \texttt{hasil.txt} menggunakan brace expansion
bersarang. Kemudian redirect daftar semua file yang dibuat ke file
\texttt{daftar-proyek.txt} menggunakan perintah \cmd{find} dan
redirection yang dipelajari di Bab~3. Tampilkan isi
\texttt{daftar-proyek.txt} dan hitung jumlah file yang berhasil dibuat.
\end{challengebox}

%% ------------------------------------------------------------------------
\section{Quoting dan Escaping}
\label{sec:quoting-escaping}

Quoting adalah salah satu konsep yang paling sering menjadi sumber
kebingungan bagi pemula. Bash melakukan banyak pemrosesan pada teks sebelum
menjalankan perintah: memecah kata berdasarkan spasi, mengembangkan
wildcard, mengganti variabel, dan menjalankan command substitution. Quoting
mengontrol pemrosesan mana yang boleh terjadi.

\textit{Single quote} (\texttt{'...'}) adalah yang paling ketat. Semua
karakter di dalamnya diperlakukan sebagai teks literal, tidak ada ekspansi
apapun yang terjadi. Bahkan tanda dolar, tanda bintang, dan backslash
semuanya dicetak apa adanya. Satu-satunya yang tidak bisa ada di dalam
single quote adalah single quote itu sendiri.

\textit{Double quote} (\texttt{"..."}) lebih longgar. Di dalamnya, Bash
masih mengembangkan variabel (\texttt{\$NAMA}), melakukan command
substitution (\texttt{\$(perintah)}), dan memproses beberapa karakter
escape backslash. Tetapi wildcard tidak dikembangkan dan spasi tidak
memecah kata. Double quote sangat berguna untuk menangani nilai variabel
yang mungkin mengandung spasi.

\textit{Backslash} (\texttt{\textbackslash}) digunakan untuk meng-escape
satu karakter berikutnya, menjadikannya literal. Misalnya
\texttt{\textbackslash\$} mencetak tanda dolar literal, dan
\texttt{\textbackslash\ } (backslash spasi) menyertakan spasi sebagai
bagian dari satu kata.

\textit{Heredoc} adalah cara menulis teks panjang tanpa perlu mengutip
setiap baris. Sintaksnya \texttt{<<EOF ... EOF} atau
\texttt{<<'EOF' ... EOF}. Dengan tanda kutip pada penanda (\texttt{'EOF'}),
isi heredoc diperlakukan seperti single quote tanpa ekspansi. Tanpa tanda
kutip, variabel dan command substitution tetap diproses.

Aturan paling penting: selalu kutip variabel yang merepresentasikan path
atau nama file dengan double quote, ditulis \texttt{"\$VARIABEL"}, bukan
\texttt{\$VARIABEL}. Jika nilai variabel mengandung spasi, Bash tanpa
kutipan akan memecahnya menjadi argumen terpisah yang bisa menyebabkan
error tak terduga.

\subsection*{Praktek 6.5: Pahami Perbedaan Quoting}

Tujuan praktek ini adalah memahami kapan menggunakan single quote,
double quote, dan backslash.

\begin{enumerate}
    \item Bandingkan perilaku single quote dan double quote:
\begin{lstlisting}[language=bash, caption={Perbedaan single quote dan double quote}]
echo '$USER bekerja di $HOME'
echo "$USER bekerja di $HOME"
echo '$USER bekerja di $HOME pada $(date)'
echo "$USER bekerja di $HOME pada $(date)"
\end{lstlisting}
    \textit{Amati:} dengan single quote, semua teks dicetak persis apa
    adanya. Dengan double quote, \var{USER}, \var{HOME}, dan hasil
    \cmd{date} semuanya diganti dengan nilai aslinya.

    \item Buat file dengan nama yang mengandung spasi, lalu akses dengan
          berbagai cara:
\begin{lstlisting}[language=bash, caption={Menangani nama file dengan spasi}]
cd ~/lab-os/chapter06-bash

touch "laporan bulan ini.txt"
touch laporan\ bulan\ lalu.txt

ls -lah
\end{lstlisting}

    \item Salin file bernam spasi menggunakan kutipan yang benar:
\begin{lstlisting}[language=bash, caption={Menyalin file dengan nama spasi secara aman}]
cp "laporan bulan ini.txt" "laporan bulan ini - salinan.txt"
cp laporan\ bulan\ lalu.txt laporan-bulan-lalu-bersih.txt

ls -lah *.txt
\end{lstlisting}

    \item Gunakan variabel untuk menyimpan path yang mengandung spasi:
\begin{lstlisting}[language=bash, caption={Variabel untuk path dengan spasi}]
FILE_LAMA="$HOME/lab-os/chapter06-bash/laporan bulan ini.txt"
FILE_BARU="$HOME/lab-os/chapter06-bash/arsip/laporan-baru.txt"

mkdir -p ~/lab-os/chapter06-bash/arsip
cp "$FILE_LAMA" "$FILE_BARU"
ls -lah ~/lab-os/chapter06-bash/arsip/
\end{lstlisting}
    \textit{Amati:} kutipan ganda pada \texttt{"\$FILE\_LAMA"} memastikan
    path yang mengandung spasi diperlakukan sebagai satu argumen, bukan
    dua argumen terpisah.

    \item Tampilkan teks yang mengandung tanda dolar literal menggunakan
          berbagai cara:
\begin{lstlisting}[language=bash, caption={Menampilkan karakter khusus secara literal}]
echo 'Harga: $50.000'
echo "Variabel PATH dimulai dengan: \$PATH"
echo "Nilai PATH sebenarnya: $PATH" | cut -c1-50
\end{lstlisting}
\end{enumerate}

\begin{challengebox}
Tulis perintah \cmd{echo} yang menampilkan baris berikut persis apa
adanya ke layar, termasuk semua karakter khusus:
\texttt{Path: \$HOME/bin/*.sh, harga: \$25, waktu: \$(date)}.
Coba dengan single quote, double quote, dan kombinasi backslash. Kemudian
buat file bernama \texttt{log [April] 2026.txt} yang berisi teks tersebut
menggunakan heredoc. Verifikasi isi file dengan \cmd{cat} dan pastikan
bisa menyalinnya ke nama file yang bersih tanpa spasi menggunakan
variabel yang dikutip dengan benar.
\end{challengebox}

%% ------------------------------------------------------------------------
\section{Rangkuman}
\label{sec:rangkuman-bash-shell}

Dalam bab ini, kamu telah mempelajari:

\begin{itemize}[noitemsep,topsep=2pt]
    \item \textbf{Shell dan konteks Bash}: Bash adalah shell default Linux
          yang membaca \file{~/.bashrc} untuk shell interaktif dan
          \file{~/.bash\_profile} untuk shell login. Perintah \cmd{source}
          menerapkan perubahan konfigurasi tanpa menutup terminal.
    \item \textbf{Prompt PS1}: variabel \var{PS1} mengontrol tampilan
          prompt. Kode khusus seperti \texttt{\textbackslash u},
          \texttt{\textbackslash h}, \texttt{\textbackslash w}, dan
          \texttt{\textbackslash t} menampilkan informasi dinamis.
    \item \textbf{Variabel dan PATH}: variabel lokal tidak diwariskan ke
          subshell, sedangkan variabel lingkungan yang dibuat dengan
          \cmd{export} diwariskan. \var{PATH} menentukan direktori tempat
          Bash mencari perintah.
    \item \textbf{Alias}: pintasan untuk perintah yang sering dipakai,
          didefinisikan dengan \texttt{alias nama='perintah'} dan
          disimpan permanen di \file{~/.bashrc}.
    \item \textbf{Fungsi shell}: lebih fleksibel dari alias karena bisa
          menerima argumen, menggunakan variabel lokal dengan \cmd{local},
          dan menjalankan logika bercabang.
    \item \textbf{Wildcard dan ekspansi}: \texttt{*}, \texttt{?}, dan
          \texttt{[...]} mencocokkan pola nama file. Brace expansion
          \texttt{\{a,b\}} dan \texttt{\{01..05\}} membuat banyak item
          sekaligus tanpa bergantung pada file yang ada.
    \item \textbf{History}: Bash menyimpan riwayat perintah yang bisa
          ditelusuri dengan panah atas, \texttt{!n}, \texttt{!!}, atau
          pencarian interaktif \keystroke{Ctrl+R}.
    \item \textbf{Quoting}: single quote membekukan semua ekspansi, double
          quote mengizinkan ekspansi variabel dan command substitution,
          backslash meng-escape satu karakter. Selalu kutip variabel path
          dengan double quote.
\end{itemize}

%% ------------------------------------------------------------------------
\section{Latihan}
\label{sec:latihan-bash-shell}

\textbf{Instruksi Umum}: kerjakan seluruh latihan di direktori
\dir{~/lab-os/chapter06-bash}. Simpan semua skrip di subdirektori
\dir{bin} dan semua hasil di subdirektori \dir{output}.

\subsubsection*{Latihan 6.1 Toolkit Bash Pribadi}
Bangun toolkit Bash pribadimu yang aktif setiap kali terminal dibuka.
Tambahkan ke \file{~/.bashrc}: minimal tiga alias yang membantu pekerjaan
harian (bukan alias yang sudah ada di sistem), satu fungsi shell
\cmd{mkcd()} yang membuat direktori dan langsung berpindah ke dalamnya,
dan satu fungsi \cmd{ringkas\_sistem()} yang menampilkan user, hostname,
uptime, dan penggunaan disk root dalam satu tampilan rapi. Tambahkan
direktori \dir{~/lab-os/chapter06-bash/bin} ke \var{PATH}. Buat skrip
\cmd{cek-toolkit} di direktori tersebut yang memanggil
\cmd{ringkas\_sistem()}. Uji bahwa \cmd{cek-toolkit} bisa dipanggil dari
\dir{/tmp} tanpa path lengkap. Simpan output pengujian ke
\file{output/hasil-toolkit.txt}.


\subsubsection*{Latihan 6.2 Audit File dengan Wildcard dan Redirection}
Buat skrip \cmd{audit-log} yang melakukan hal berikut. Pertama, cari
semua file \texttt{*.log} di dalam \dir{/var/log} menggunakan wildcard
dan simpan daftarnya ke \file{output/daftar-log.txt}, sedangkan pesan
error dari file yang tidak bisa dibaca disimpan ke \file{output/error-log.txt}
secara terpisah. Kedua, buat 10 file log contoh di
\dir{~/lab-os/chapter06-bash/logs} menggunakan brace expansion dengan nama
\texttt{app-\{01..05\}.log} dan \texttt{svc-\{01..05\}.log}. Ketiga,
gunakan wildcard untuk memindahkan semua file \texttt{svc-*.log} ke
subdirektori \dir{arsip/} setelah melakukan preview dengan \cmd{echo}
terlebih dahulu. Keempat, buat arsip \texttt{tar.gz} dari semua log yang
tersisa. Simpan seluruh perintah yang digunakan ke
\file{output/riwayat-audit.txt} menggunakan \cmd{history}.


\subsubsection*{Latihan 6.3 Sistem Health Check Harian}
Buat skrip \cmd{health-check} di direktori \dir{~/bin} yang menampilkan
ringkasan kondisi sistem. Skrip harus menggunakan minimal dua alias
(yang kamu definisikan sendiri), satu fungsi shell dengan variabel lokal,
variabel lingkungan kustom untuk path output, dan perintah \cmd{tee}
untuk menyimpan output ke file log sekaligus menampilkannya di terminal.
Output skrip mencakup: tanggal dan waktu, nama pengguna dan host, uptime,
penggunaan memori, penggunaan disk setiap filesystem. Simpan log ke
\file{output/health-TANGGAL.log} dengan tanggal otomatis dari
command substitution. Tambahkan penanganan sederhana: jika penggunaan
disk lebih dari 80\%, tampilkan teks ``PERINGATAN'' berwarna merah
menggunakan kode ANSI.


%% ------------------------------------------------------------------------
\section{Referensi}
\label{sec:referensi-bash-shell}

\begin{itemize}
    \item Free Software Foundation. \textit{GNU Bash Reference Manual}.
          \url{https://www.gnu.org/software/bash/manual/}
    \item Shotts, William. \textit{The Linux Command Line}, 2nd Edition.
          No Starch Press, 2019.
    \item Robbins, Arnold dan Beebe, Nelson H.F. \textit{Classic Shell
          Scripting}. O'Reilly Media, 2005.
    \item Blum, Richard dan Bresnahan, Christine. \textit{Linux Command
          Line and Shell Scripting Bible}, 4th Edition. Wiley, 2021.
    \item Nemeth, Evi dkk. \textit{UNIX and Linux System Administration
          Handbook}, 5th Edition. Pearson, 2017.
\end{itemize}
